\documentclass[10pt,a4paper]{amsart}
\usepackage[margin=19mm]{geometry}
\usepackage{amsmath,amssymb,mathtools}
\usepackage[scr=rsfs,cal=euler]{mathalfa}
\usepackage{xfrac}
\usepackage[unicode,hidelinks,bookmarks=false]{hyperref}
\newcommand{\ie}{i.e.\ }
\newcommand{\cf}{cf.\ }
\newcommand{\resp}{respectively\ }
\newcommand{\Subsection}{\subsection}
\providecommand{\nobreakdash}{\nobreak\hbox{-}}
\setlength{\emergencystretch}{2em}
\multlinegap0pt
\usepackage{amssymb}
\usepackage{xcolor}
\usepackage{cancel}
\newtheorem{rem}{Remark}
\newtheorem{thm}{Theorem}
\newcommand{\CE}{\mathcal{E}}
\newcommand{\CM}{\mathcal{M}}
\newcommand{\CN}{\mathcal{N}}
\newcommand{\CP}{\mathcal{P}}
\newcommand{\auleq}{\underset{a.u.}{\leq }}
\newcommand{\norm}[1]{\left\Vert#1\right\Vert}
\title{Non-commutative Law of iterated logarithm}
\author{Sourav Panja, \'Eric Ricard, Diptesh Saha}
\date{Source: arXiv:2509.22037v2 (CC BY 4.0)}
\begin{document}
\maketitle

\noindent\emph{Source and licence.} Non-commutative Law of iterated logarithm. Version arXiv:2509.22037v2 (CC BY 4.0), distributed by arXiv under the Creative Commons Attribution 4.0 International licence, http://creativecommons.org/licenses/by/4.0/, as recorded in the arXiv licence metadata for this exact version and shown on the abstract page https://arxiv.org/abs/2509.22037v2. Full licence notice: \texttt{licenses/arxiv-2509.22037-CC-BY-4.0.txt}.\par\smallskip

\noindent\textit{Editorial scope.} Counted unit: the article's thm trick (source0.tex lines 872--877). The article's complete original proof of it is delivered with it (source0.tex lines 878--936, one \texttt{proof} environment of 59 lines). No mathematics is written, repaired or re-derived: the statement and the proof are the article's own lines, in the article's own notation, and no retained line is substituted or re-flowed.  Retained uncounted as the source-local context the proof needs: lines 423--446; lines 720--725.  Nothing was omitted from any retained span, so the package carries no omission record.  No retained line contains a citation, so the wrapper carries no bibliography and no bibliography entry.  No retained line contains a figure environment, an image-inclusion command or drawing code, so no image is delivered and the package declares no asset dependency.  Editorial formatting only, in addition: an AMS \texttt{amsart} preamble that restates the article's own declarations and macros used by the retained lines, namely the environments \texttt{rem}, \texttt{thm} on separate counters, together with the standard packages those lines need.  The retained environments print in this document as Remark 1, Theorem 1 in order of appearance, whereas the article prints the same items with the numbering of its own full document.  The complete native source of the article as distributed in the arXiv e-print for this exact version is bundled unchanged as \texttt{sources/185786/source0.tex}.
\begin{rem}\label{rem: indep of ambient vna}
Let $\CN$ be a von Neumann subalgebra of $\CM$. Let $(x_n)$ be a sequence in $\CN$ with $\limsup_{n \to \infty} x_n \auleq K $ in $\CM$ for some constant $K>0$. Let $\epsilon, \delta>0$. We first choose $0< \eta <1$ such that $(K+ \frac{\delta}{2})/ (1-\eta) \leq K+ \delta$. Then, there exists $p \in \CP(\CM)$ such that $\tau(1-p)< \eta \epsilon$ and 
\[
\limsup_{n \to \infty} \norm{x_np} \leq K + \frac{\delta}{2}.
\]
We claim that there exists $q \in \CP(\CN)$ such that $\tau(1-q)< \epsilon$ and 
\[
\limsup_{n \to \infty} \norm{x_nq} \leq K + \delta.
\]


Consider $p':= E(p)$, where $E$ is the faithful, normal conditional expectation from $\CM$ onto $\CN$. Now, consider $q:= 1_{[1-\eta, 1]}(p')\in \CN$. Next, observe that $1- p' \geq \eta (1-q)$ and
\[
\eta \epsilon \geq \tau(1-p)= \tau(1-p') \geq \eta \tau(1-q).
\]
On the other hand, since $E(p)\geq (1-\eta) q$, we have \[
(1-\eta)\norm{x_nq}^2=(1-\eta)\norm{x_nqx_n^*} \leq  \norm{x_n E(p) x_n^*}=\norm{E(x_n p x_n^*)} \leq  \norm{x_np}^2.
\]
Therefore, we obtain $\tau(1-q) \leq \epsilon$ and
\[
\limsup_{n \to \infty} \norm{x_nq} \leq \frac{1}{1-\eta} \limsup_{n \to \infty} \norm{x_np} \leq K+ \delta.
\]
A similar argument also holds true for almost uniform convergence.
\end{rem}

\begin{thm}\label{thm: main thm}
    Let $0=x_0, x_1, x_2, \ldots$ be a sequence of self-adjoint martingales in the von Neumann algebra $(\mathcal{M}, \tau)$. Let $s_n^2=\norm{\sum_{k=1}^n E_{k-1}(d_k^2)}$ and $u_n=L(s_n^2)^{1/2}$. Suppose that $s_n^2 \to \infty$ and $\norm{d_n} \leq \frac{\alpha_n s_n}{u_n}$, where $(\alpha_n)$ is a sequence of positive {\color{blue} real} numbers such that $\alpha_n \to 0$, then 
    \[
    \limsup_{n \to \infty} \frac{x_n}{s_n u_n} \underset{a.u.}{\leq }\sqrt{2}.
    \]
\end{thm}

\begin{thm}\label{trick}
    Let $(x_n)_{n \geq 0}$ be a non-commutative martingale in $(\mathcal{M}, \tau)$ with $x_0 = 0$. Define \newline $t_n^2 = \max\left\{\norm{\sum_{i=1}^n E_{i-1}(d_i^*d_i)},\norm{\sum_{i=1}^n E_{i-1}(d_id_i^*)}\right\}$ and $v_n = [L(t_n^2)]^{1/2}$. Suppose that $t_n^2 \to \infty$ and $\norm{d_n}_\infty \leq \frac{\alpha_n t_n}{v_n}$, where $(\alpha_n)$ is a sequence of positive numbers such that $\alpha_n \to 0$, then 
    $$
    \limsup_{n \to \infty} \frac{x_n}{t_n v_n} \underset{a.u.}{\leq }\sqrt{2}.
    $$
\end{thm}

\begin{proof}
    Consider the von Neumann algebras 
    $$
    \tilde{\mathcal{M}} := \mathcal{M} \otimes M_2(\mathbb{C}) \quad \text{and} \quad \tilde{\mathcal{M}}_k := \mathcal{M}_k \otimes M_2(\mathbb{C})
    $$
    for each integer $k \geq 1$. Let $(\tilde{x}_n)_n$ be a sequence of elements in $\tilde{\mathcal{M}}$ defined by
    $$
    \tilde{x}_n := \begin{pmatrix}
    0 & x_n \\
    x_n^* & 0
    \end{pmatrix}, \quad n \geq 0,
    $$
    where $x_n \in \mathcal{M}$.
    
    The sequence $(\tilde{x}_n)_n$ forms a self-adjoint martingale with respect to the increasing filtration $(\tilde{\mathcal{M}}_n, \mathcal{E}_n)$, where the conditional expectation 
    $\mathcal{E}_n : \tilde{\mathcal{M}}_{n+1} \to \tilde{\mathcal{M}}_n$
    satisfies
    $$
    \mathcal{E}_n(\tilde{x}_{n+1}) := \begin{pmatrix}
    0 & E_n(x_{n+1}) \\
    E_n(x_{n+1}^*) & 0
    \end{pmatrix}, \quad n \geq 0.
    $$
    Here, $E_n$ denotes the conditional expectation from $\mathcal{M}_{n+1}$ onto $\mathcal{M}_n$.
    Let us define for $n\geq1$,
    $\tilde{d}_n := \tilde{x}_n - \tilde{x}_{n-1}.$
    Observe that each $\tilde{d}_n$ is self-adjoint for all $n \geq 0$. Moreover,
    \[
        \tilde{d}_n^2 = 
        \begin{pmatrix} 0 & d_n \\ d_n^* & 0 \end{pmatrix} 
        \begin{pmatrix} 0 & d_n \\ d_n^* & 0 \end{pmatrix} =
        \begin{pmatrix} d_n d_n^* & 0 \\ 0 & d_n^* d_n \end{pmatrix}.
    \]
    Define, for each $n \geq 1$,
    \[
        \tilde{t}_n^2 := \left\| \sum_{i=1}^n \CE_{i-1}\left(\tilde{d}_i^2\right) \right\|.
    \]
    It follows that $\tilde{t}_n = t_n$ for all $n \geq 1$. Furthermore, for every $n \geq 1$,
    \[
        \|\tilde{d}_n\|_\infty = \|d_n^* d_n\|^{1/2} = \|d_n\|.
    \]
    Consequently, $(\tilde{x}_n)_n$ constitutes a sequence of self-adjoint martingales in the von Neumann algebra $\tilde{\mathcal{M}}$ such that
    \[
        \lim_{n \to \infty} \tilde{t}_n = \infty
        \quad \text{and} \quad
        \|\tilde{d}_n\|_\infty \leq \frac{\alpha_n \tilde{t}_n}{v_n}, \quad \text{for all } n \geq 1.
    \]
    Applying Theorem~\ref{thm: main thm}, we deduce that
    \[
        \limsup_{n \to \infty} \frac{\tilde{x}_n}{t_n v_n} \underset{a.u.}{\leq} \sqrt{2}.
    \]

    Taking $\tilde r=\begin{pmatrix} 0 & 0 \\ 1 & 0 \end{pmatrix}$ in $\tilde \CM$ we first note that $\tilde{r} \tilde{x}_n = \begin{pmatrix} 0 & 0 \\ 0 & x_n \end{pmatrix}$ and obtain, 
    $$\limsup_{n\to \infty} \frac{\tilde{r}\tilde{x}_n }{t_n v_n} \underset{{a.u.}}{\leq} \sqrt{2}.$$
    Since almost uniform boundedness does not depend on the ambient von Neumann algebra (cf. Remark \ref{rem: indep of ambient vna}), we obtain 
    \[
    \limsup_{n\to \infty} \frac{x_n}{t_n v_n} \underset{{a.u.}}{\leq} \sqrt{2}.
    \]
\end{proof}

\end{document}
